\documentclass[aspectratio=169]{beamer}
\usepackage[english]{babel}
\usepackage[utf8]{inputenc}
\usepackage[T1]{fontenc}
\usepackage{lmodern}
\usepackage{listings}
\usepackage{xcolor}
\usepackage{graphicx}
\usepackage{textcomp}
\DeclareUnicodeCharacter{2014}{---}
\DeclareUnicodeCharacter{2013}{--}
\DeclareUnicodeCharacter{2212}{-}
\DeclareUnicodeCharacter{2192}{\ensuremath{\rightarrow}}
\DeclareUnicodeCharacter{00B1}{\ensuremath{\pm}}
\DeclareUnicodeCharacter{00D7}{\ensuremath{\times}}
\DeclareUnicodeCharacter{00B7}{\ensuremath{\cdot}}
\DeclareUnicodeCharacter{20AC}{EUR}
\DeclareUnicodeCharacter{2070}{\textsuperscript{0}}
\DeclareUnicodeCharacter{00B9}{\textsuperscript{1}}
\DeclareUnicodeCharacter{00B2}{\textsuperscript{2}}
\DeclareUnicodeCharacter{00B3}{\textsuperscript{3}}
\DeclareUnicodeCharacter{2074}{\textsuperscript{4}}
\DeclareUnicodeCharacter{2075}{\textsuperscript{5}}
\DeclareUnicodeCharacter{2076}{\textsuperscript{6}}
\DeclareUnicodeCharacter{2077}{\textsuperscript{7}}
\DeclareUnicodeCharacter{2078}{\textsuperscript{8}}
\DeclareUnicodeCharacter{2079}{\textsuperscript{9}}
\DeclareUnicodeCharacter{2248}{\ensuremath{\approx}}
\DeclareUnicodeCharacter{2264}{\ensuremath{\le}}

% Colors for code
\definecolor{codebg}{RGB}{245,245,245}
\definecolor{codecomment}{RGB}{0,128,0}
\definecolor{codekeyword}{RGB}{0,0,255}
\definecolor{codestring}{RGB}{163,21,21}
\definecolor{bandcolor}{RGB}{200,50,50}

\lstdefinestyle{pythonstyle}{
    language=Python,
    backgroundcolor=\color{codebg},
    basicstyle=\ttfamily\scriptsize,
    keywordstyle=\color{codekeyword},
    commentstyle=\color{codecomment},
    stringstyle=\color{codestring},
    showstringspaces=false,
    breaklines=true,
    frame=single,
    rulecolor=\color{black},
    escapeinside={(*@}{@*)},
    moredelim=**[l][\color{bandcolor}\bfseries]{\#\ ---\ NEW},
    moredelim=**[l][\color{bandcolor}\bfseries]{\#\ ---\ CHANGED},
    literate=
      {á}{{\'a}}1 {é}{{\'e}}1 {í}{{\'i}}1 {ó}{{\'o}}1 {ú}{{\'u}}1
      {Á}{{\'A}}1 {É}{{\'E}}1 {Í}{{\'I}}1 {Ó}{{\'O}}1 {Ú}{{\'U}}1
      {ñ}{{\~n}}1 {Ñ}{{\~N}}1 {ü}{{\"u}}1 {Ü}{{\"U}}1
      {¿}{{?`}}1 {¡}{{!`}}1
      {—}{{---}}1 {–}{{--}}1 {−}{{-}}1
      {±}{{\ensuremath{\pm}}}1 {×}{{\ensuremath{\times}}}1
      {→}{{\ensuremath{\rightarrow}}}1
      {°}{{\textdegree}}1
      {·}{{\ensuremath{\cdot}}}1
      {€}{{\texteuro}}1
      {≈}{{\ensuremath{\approx}}}1
      {≤}{{\ensuremath{\le}}}1
      {⁰}{{\textsuperscript{0}}}1 {¹}{{\textsuperscript{1}}}1
      {²}{{\textsuperscript{2}}}1 {³}{{\textsuperscript{3}}}1
      {⁴}{{\textsuperscript{4}}}1 {⁵}{{\textsuperscript{5}}}1
      {⁶}{{\textsuperscript{6}}}1 {⁷}{{\textsuperscript{7}}}1
      {⁸}{{\textsuperscript{8}}}1 {⁹}{{\textsuperscript{9}}}1
}

\lstset{style=pythonstyle}

% Title slide macro: \lessontitle{Lesson Number}{Title}
\newcommand{\lessontitle}[2]{
    \title{Lesson #1: #2}
    \author{}
    \institute{Tec de Monterrey}
    \begin{frame}[plain]
        \titlepage
    \end{frame}
}

% Figure frame macro: \figureframe{Title}{Image path}
\newcommand{\figureframe}[2]{
    \begin{frame}{#1}
        \begin{center}
            \includegraphics[height=0.8\textheight,width=\textwidth,keepaspectratio]{#2}
        \end{center}
    \end{frame}
}


% Listings pulled from src/ with \lstinputlisting, so the slide is the file.
\lstset{
    basicstyle=\ttfamily\fontsize{8pt}{9.6pt}\selectfont,
    breaklines=true,
    breakatwhitespace=true,
    columns=fullflexible,
    keepspaces=true,
    xleftmargin=0.3em,
    framesep=2pt,
    aboveskip=0pt,
    belowskip=0pt
}

\newcommand{\resultframe}[3]{%
    \begin{frame}{#1}
        \begin{center}
            \includegraphics[height=0.62\textheight,width=\textwidth,keepaspectratio]{#2}
        \end{center}
        \vspace{0.25em}
        {\small #3\par}
    \end{frame}%
}

\newcommand{\claimframe}[2]{%
    \begin{frame}{#1}
        {\small #2\par}
    \end{frame}%
}

\date{}

\begin{document}

\lessontitle{04}{Crossover}

\section{This lesson}

\begin{frame}{This lesson}
Lesson 03 opened selection on one generation. Crossover is next, still on two fixed parents, with one instrument.

\vspace{0.6em}
\begin{itemize}
    \item A crossover is defined by the children it can reach.
    \item More cut points is not a strength dial.
    \item Cut-and-swap copies parental values; blend computes new ones --- and leaves the interval.
\end{itemize}
\end{frame}

% ================= offspring_01_two_parents =================
\begin{frame}{Reach is the support of an operator}
For BLX-\(\alpha\), let \(m=\min(a,b)\), \(M=\max(a,b)\), and \(d=M-m\):
\[
x'\sim U[m-\alpha d,\,M+\alpha d],\qquad
\Pr(x'\notin[m,M])=\frac{2\alpha}{1+2\alpha}.
\]
Uniform crossover has \(2^L\) possible first children when all \(L\) parental
gene pairs differ; linear crossover reaches two affine children. Canonical
(n)-point cuts use (1,\ldots,L-1), but this code excludes (L-1).
Samples illustrate reach; enumeration establishes exact cardinality. OX1
preserves permutation membership and uniqueness.
\end{frame}

\section{Two parents, and how to measure what an operator does to them}

\begin{frame}{Two parents, and how to measure what an operator does to them}
Crossover is the operator that takes two chromosomes and returns two new ones.
Nothing else in a genetic algorithm creates a combination that was not there
before, so the whole lesson turns on one question: GIVEN THESE TWO PARENTS,
WHAT IS THE SET OF CHILDREN AN OPERATOR CAN PRODUCE?
\end{frame}

\resultframe{Two parents, and how to measure what an operator does to them}{../figures/offspring_01_two_parents.png}{%
    The zero row: 2 children reached, 0.00\% new genes, 100.00\% clones, 100.00\% legal. Also the number every later step is measured against: choosing freely between the parents allows 2⁶ = \textbf{64} chromosomes}

% ================= offspring_02_one_point =================
\section{One-point crossover, and the size of its reach}

\begin{frame}{One-point crossover, and the size of its reach}
The textbook operator: cut both parents at the same point, swap the tails. It is
worth starting here not because it is the best but because its reachable set is
small enough to write down in full - which is the only way to see what the
operator can and cannot do.
\end{frame}

\begin{frame}[fragile]{(1) crossover\_one\_point()}
\lstinputlisting[basicstyle=\ttfamily\fontsize{8pt}{9.6pt}\selectfont,firstline=129,lastline=142]{../src/offspring_02_one_point.py}
\end{frame}

\begin{frame}[fragile]{(2) the reachable set}
\lstinputlisting[basicstyle=\ttfamily\fontsize{6pt}{7.2pt}\selectfont,firstline=150,lastline=182]{../src/offspring_02_one_point.py}
\end{frame}

\begin{frame}[fragile]{(3) the cut figure}
\lstinputlisting[basicstyle=\ttfamily\fontsize{7pt}{8.5pt}\selectfont,firstline=185,lastline=210]{../src/offspring_02_one_point.py}
\end{frame}

\resultframe{One-point crossover, and the size of its reach}{../figures/offspring_02_one_point.png}{%
    2,000 draws produce \textbf{10} distinct children; enumerating the 5 cut points produces the same 10. \textbf{0.00\%} new genes: the operator chooses which parent, never what value}

\resultframe{The mechanism: one cut, two children}{../figures/offspring_02_one_point_schematic.png}{%
    Cut after gene 3. Blue genes come from parent 1, red genes from parent 2.
    The child is one parent up to the cut and the other after it. No gene is invented.}

% ================= offspring_03_more_cuts =================
\section{More cut points, and what the family can never do}

\begin{frame}{More cut points, and what the family can never do}
If one cut is good, are three better? Measured on this pair the answer is no,
and the reason is worth more than the answer. What separates the members of this
family is not how many cuts they make but how many of the 64 combinations of the
parents' genes they can actually reach: uniform crossover reaches all of them,
the n-point operators reach a handful.
\end{frame}

\begin{frame}[fragile]{(1) crossover\_n\_point()}
\lstinputlisting[basicstyle=\ttfamily\fontsize{7pt}{8.5pt}\selectfont,firstline=148,lastline=168]{../src/offspring_03_more_cuts.py}
\end{frame}

\begin{frame}[fragile]{(2) crossover\_uniform()}
\lstinputlisting[basicstyle=\ttfamily\fontsize{8pt}{9.6pt}\selectfont,firstline=171,lastline=188]{../src/offspring_03_more_cuts.py}
\end{frame}

\begin{frame}[fragile]{(3) the reach table}
\lstinputlisting[basicstyle=\ttfamily\fontsize{6pt}{7.2pt}\selectfont,firstline=196,lastline=243]{../src/offspring_03_more_cuts.py}
\end{frame}

\begin{frame}[fragile]{(4) the pattern figure}
\lstinputlisting[basicstyle=\ttfamily\fontsize{7pt}{8.5pt}\selectfont,firstline=246,lastline=268]{../src/offspring_03_more_cuts.py}
\end{frame}

\resultframe{More cut points, and what the family can never do}{../figures/offspring_03_more_cuts.png}{%
    Reach is not monotone in cuts: one-point \textbf{10}, 2-point \textbf{12}, 3-point \textbf{8}, uniform \textbf{64 of 64}. Uniform's price: \textbf{3.20\%} of its children are a parent handed back}

% ================= offspring_04_blend =================
\section{Arithmetic on genes, and the first illegal children}

\begin{frame}{Arithmetic on genes, and the first illegal children}
Every operator so far chose between the parents' genes. These two compute new
ones - which is only possible because a gene here is a NUMBER, and numbers can
be averaged, interpolated and extrapolated. That single assumption about the
representation is what step 6 will take away.
\end{frame}

\begin{frame}[fragile]{(1) clamp()}
\lstinputlisting[basicstyle=\ttfamily\fontsize{8pt}{9.6pt}\selectfont,firstline=191,lastline=200]{../src/offspring_04_blend.py}
\end{frame}

\begin{frame}[fragile]{(2) crossover\_blend()}
\lstinputlisting[basicstyle=\ttfamily\fontsize{6pt}{7.2pt}\selectfont,firstline=203,lastline=233]{../src/offspring_04_blend.py}
\end{frame}

\begin{frame}[fragile]{(3) crossover\_linear()}
\lstinputlisting[basicstyle=\ttfamily\fontsize{7pt}{8.5pt}\selectfont,firstline=236,lastline=254]{../src/offspring_04_blend.py}
\end{frame}

\begin{frame}[fragile]{(4) the reach table}
\lstinputlisting[basicstyle=\ttfamily\fontsize{6pt}{7.2pt}\selectfont,firstline=262,lastline=348]{../src/offspring_04_blend.py}
\end{frame}

\begin{frame}[fragile]{(5) the blend figure}
\lstinputlisting[basicstyle=\ttfamily\fontsize{6pt}{7.2pt}\selectfont,firstline=351,lastline=380]{../src/offspring_04_blend.py}
\end{frame}

\resultframe{Arithmetic on genes, and the first illegal children}{../figures/offspring_04_blend.png}{%
    The first computed genes (\textbf{99.57\%} new) and the first illegal children. alpha = 0: \textbf{0.00\%} of genes beyond the parents, 100\% legal — Lesson 01's inward collapse, measured. alpha = 0.5: \textbf{49.63\%} beyond, only \textbf{50.08\%} legal; with \texttt{clamp()}, 100\% legal and still 49.63\% beyond. Linear: 100.00\% new genes and exactly \textbf{2} reachable children}

% ================= offspring_05_fitness_driven =================
\section{Letting the operator look at what it made}

\begin{frame}{Letting the operator look at what it made}
Every operator so far was blind: it produced two children and handed them over
without ever asking whether they were any good. Gridin's fitness-driven variant
adds one line - rank the two parents and the two children together, return the
best two - and with it the guarantee that crossover can never make the pair
worse.
\end{frame}

\begin{frame}[fragile]{(1) objective()}
\lstinputlisting[basicstyle=\ttfamily\fontsize{8pt}{9.6pt}\selectfont,firstline=252,lastline=261]{../src/offspring_05_fitness_driven.py}
\end{frame}

\begin{frame}[fragile]{(2) crossover\_fitness\_driven()}
\lstinputlisting[basicstyle=\ttfamily\fontsize{8pt}{9.6pt}\selectfont,firstline=264,lastline=280]{../src/offspring_05_fitness_driven.py}
\end{frame}

\begin{frame}[fragile]{(3) the survivors}
\lstinputlisting[basicstyle=\ttfamily\fontsize{6pt}{7.2pt}\selectfont,firstline=288,lastline=357]{../src/offspring_05_fitness_driven.py}
\end{frame}

\begin{frame}[fragile]{(4) the fitness figure}
\lstinputlisting[basicstyle=\ttfamily\fontsize{6pt}{7.2pt}\selectfont,firstline=360,lastline=391]{../src/offspring_05_fitness_driven.py}
\end{frame}

\resultframe{Letting the operator look at what it made}{../figures/offspring_05_fitness_driven.png}{%
    Blind blend is worse than its parents in \textbf{41.70\%} of draws — crossover guarantees nothing. The fitness-driven rule is worse in \textbf{0.00\%}, and pays for it: \textbf{16.50\%} of draws return the parents unchanged, and only \textbf{56.10\%} of what it returns is a child}

% ================= offspring_06_the_divide =================
\section{The same operators, on a chromosome that is a permutation}

\begin{frame}{The same operators, on a chromosome that is a permutation}
Nothing about the operators changes in this step. The chromosome changes: nine
genes instead of six, and the genes are the nine stops of a delivery round, each
of which must appear exactly once.
\end{frame}

\begin{frame}[fragile]{(1) the nine stops}
\lstinputlisting[basicstyle=\ttfamily\fontsize{7pt}{8.5pt}\selectfont,firstline=279,lastline=297]{../src/offspring_06_the_divide.py}
\end{frame}

\begin{frame}[fragile]{(2) tour\_length()}
\lstinputlisting[basicstyle=\ttfamily\fontsize{8pt}{9.6pt}\selectfont,firstline=300,lastline=305]{../src/offspring_06_the_divide.py}
\end{frame}

\begin{frame}[fragile]{(3) is\_legal\_tour()}
\lstinputlisting[basicstyle=\ttfamily\fontsize{8pt}{9.6pt}\selectfont,firstline=308,lastline=328]{../src/offspring_06_the_divide.py}
\end{frame}

\begin{frame}[fragile]{(4) the same operators}
\lstinputlisting[basicstyle=\ttfamily\fontsize{6pt}{7.2pt}\selectfont,firstline=336,lastline=408]{../src/offspring_06_the_divide.py}
\end{frame}

\begin{frame}[fragile]{(5) the map figure}
\lstinputlisting[basicstyle=\ttfamily\fontsize{6pt}{7.2pt}\selectfont,firstline=411,lastline=444]{../src/offspring_06_the_divide.py}
\end{frame}

\resultframe{The same operators, on a chromosome that is a permutation}{../figures/offspring_06_the_divide.png}{%
    The same operators, a route instead of a point. \textbf{8 of 8} cut points give a route that visits one stop twice and another never. Legal children: one-point \textbf{0.00\%}, 3-point \textbf{0.00\%}, blend \textbf{0.00\%}, linear \textbf{0.00\%}, uniform \textbf{0.95\%} — and the count of uniform's legal children that were not simply a parent is \textbf{0}}

% ================= offspring_07_ordered =================
\section{Ordered crossover, and the divide read in both directions}

\begin{frame}{Ordered crossover, and the divide read in both directions}
Ordered crossover (OX1) is the answer to step 6: it never swaps gene values
position by position, so a permutation goes in and a permutation comes out, by
construction rather than by luck. This step checks that it is legal, checks what
it actually inherits, and then does the thing that closes the lesson - hands OX1
the real-valued parents of step 1 and watches what happens.
\end{frame}

\begin{frame}[fragile]{(1) crossover\_order()}
\lstinputlisting[basicstyle=\ttfamily\fontsize{6pt}{7.2pt}\selectfont,firstline=327,lastline=358]{../src/offspring_07_ordered.py}
\end{frame}

\begin{frame}[fragile]{(2) the adjacency test}
\lstinputlisting[basicstyle=\ttfamily\fontsize{6pt}{7.2pt}\selectfont,firstline=366,lastline=420]{../src/offspring_07_ordered.py}
\end{frame}

\begin{frame}[fragile]{(3) the other direction}
\lstinputlisting[basicstyle=\ttfamily\fontsize{7pt}{8.5pt}\selectfont,firstline=423,lastline=446]{../src/offspring_07_ordered.py}
\end{frame}

\begin{frame}[fragile]{(4) the whole comparison}
\lstinputlisting[basicstyle=\ttfamily\fontsize{6pt}{7.2pt}\selectfont,firstline=449,lastline=517]{../src/offspring_07_ordered.py}
\end{frame}

\begin{frame}[fragile]{(5) the divide figure}
\lstinputlisting[basicstyle=\ttfamily\fontsize{7pt}{8.5pt}\selectfont,firstline=520,lastline=541]{../src/offspring_07_ordered.py}
\end{frame}

\resultframe{Ordered crossover, and the divide read in both directions}{../figures/offspring_07_ordered.png}{%
    OX1 is legal \textbf{100.00\%} on routes and keeps \textbf{7.17 of 9} parent legs, against \textbf{3.71} for a random route. Handed the real-valued parents it does not crash: it reaches \textbf{2} children, \textbf{100.00\%} of them a parent. Final table: operators useful on both representations — \textbf{0}}

\section{Next}

\begin{frame}{Next}
On a route of nine stops, one-point, n-point, blend and linear produce a legal child 0\% of the time. Ordered crossover is legal 100\% on routes and silently returns the parents on real genes. Zero operators work on both.

\vspace{0.8em}
\textbf{Lesson 05 --- Mutation}

The third operator. OX1 on a tour was a silent no-op here; that bill is paid in Lesson 10, where order is the chromosome.
\end{frame}
\end{document}
